\subsection{Aut(g) 上的 Zariski 拓扑}

命题 8. 设 V 是向量空间 g 的自同态组成的集合。则 \(\operatorname{Aut}(\mathfrak{g})\) 在 V 中对 Zariski 拓扑是闭的（第七章附录 I）。

设 \(K\) 是 \(\mathfrak{g}\) 的 Killing 型。若 \(s \in \operatorname{Aut}(\mathfrak{g})\) ，

\[
  [ s x, s y ] = [ x, y ]\tag{2}
\]

\[
\mathrm{K} (s x, s y) = \mathrm{K} (x, y)\tag{3}
\]

对一切 \(x, y \in g\) 成立。反之，设 s 是 V 的元素，满足对一切 \(x, y \in g\) 的 (2) 与 (3)。则 \(\operatorname{Ker}(s) = 0\) ，故 s 是双射且 \(s \in \operatorname{Aut}(\mathfrak{g})\) 。但是，对一切 \(x, y \in g\) ，从 V 到 g 与 k 的映射 \(s \mapsto [sx, sy]\) 与 \(s \mapsto \operatorname{K}(sx, sy)\) 是多项式的。

命题 9. 设 \(\mathfrak{h}\) 是 \(\mathfrak{g}\) 的一个分裂嘉当子代数。

\begin{enumerate}
\def\labelenumi{(\roman{enumi})}
\item
  群 \(f(\mathrm{T}_{\mathrm{Q}})\) 在 \(\operatorname{Aut}(\mathfrak{g})\) 中对 Zariski 拓扑是闭的。
\item
  群 \(f(q(\mathrm{T_P}))\) 在 \(f(\mathrm{T_Q})\) 中对 Zariski 拓扑是稠密的。
\end{enumerate}

断言 (i) 由等式 \(f(\mathrm{T}_{\mathrm{Q}}) = \operatorname{Aut}(\mathfrak{g},\mathfrak{h})\cap \operatorname{Ker}\varepsilon\) （命题 2）得出。置 \(m = (\mathrm{P}(\mathrm{R}):\mathrm{Q}(\mathrm{R}))\) 。设 \(\mathrm{F}\) 是 \(\mathrm{V}\) 上的一个多项式函数；我们假设 \(\mathrm{F}\) 在每个元素的 \(m\) 次幂（元素属于 \(f(\mathrm{T}_{\mathrm{Q}})\) ）上为零，并证明 \(\mathrm{F}|f(\mathrm{T}_{\mathrm{Q}}) = 0\) ；鉴于引理 3，这就将证明 (ii)。

集合 \(V'\) （V 中在 h 上诱导恒等且使每个 \(g^{\alpha}\) 稳定的元素组成的集合）可以与 \(k^{R}\) 等同。设 \(F'\) 是 F 在 \(V' = k^{R}\) 上的限制；这是一个多项式函数。我们有 \(f(\mathrm{T}_{\mathrm{Q}}) \subset \mathrm{V}'\) 。设 \(\mathrm{B} = (\alpha_{1}, \ldots, \alpha_{l})\) 是 R 的一个基。对一切 \(t = (t_{1}, \ldots, t_{l}) \in k^{*B}\) ，用 \(\varphi(t)\) 表示从 \(\mathrm{Q}(\mathrm{R})\) 到群 \(k^{*}\) 的扩张 t 的同态。则 \(\mathrm{F}'(f(\varphi(t)))\) 可以写成有限和

\[
\sum_ {n _ {1}, \dots , n _ {l} \in \mathbf {Z}} c _ {n _ {1}, \dots , n _ {l}} t _ {1} ^ {n _ {1}} \dots t _ {l} ^ {n _ {l}} = \mathrm{H} (t _ {1}, \dots , t _ {l}).
\]

由假设，

\[
0 = \mathrm{H} (t _ {1} ^ {m}, \dots , t _ {l} ^ {m}) = \sum_ {n _ {1}, \dots , n _ {l} \in \mathbf {Z}} c _ {n _ {1}, \dots , n _ {l}} t _ {1} ^ {m n _ {1}} \dots t _ {l} ^ {m n _ {l}}
\]

对一切 \(t_{1},\ldots,t_{l}\in k^{*}\) 成立。于是诸 \(c_{n_{1},\ldots,n_{l}}\) 是一个在 \(k^{*l}\) 上为零的 l 个变量的多项式的系数；因此它们全为零。

命题 10. 假设 \(\mathfrak{g}\) 是可分裂的。

\begin{enumerate}
\def\labelenumi{(\roman{enumi})}
\item
  群 \(\operatorname{Aut}_{e}(\mathfrak{g})\) 在 \(\operatorname{Aut}_{0}(\mathfrak{g})\) 中对 Zariski 拓扑是稠密的。
\item
  群 \(\mathrm{Aut}_e(\mathfrak{g})\) 与 \(\mathrm{Aut}_0(\mathfrak{g})\) 对 Zariski 拓扑是连通的。
\end{enumerate}

由命题 3，\(f(q(\mathrm{T}_{\mathrm{P}})) \subset \mathrm{Aut}_e(\mathfrak{g})\) 。对一切 \(s \in \mathrm{Aut}_e(\mathfrak{g})\) ，\(s.f(q(\mathrm{T}_{\mathrm{P}}))\) 在 Zariski 拓扑中的闭包由命题 9 知包含 \(s.f(\mathrm{T}_{\mathrm{Q}})\) 。因此 \(\mathrm{Aut}_e(\mathfrak{g})\) 的闭包包含 \(\mathrm{Aut}_e(\mathfrak{g}).f(\mathrm{T}_{\mathrm{Q}}) = \mathrm{Aut}_0(\mathfrak{g})\) （命题 6）。这就证明了 (i)。

设 \(\mathrm{Aut}_e(\mathfrak{g}) = \Omega \cup \Omega'\) 是 \(\mathrm{Aut}_e(\mathfrak{g})\) 的一个分割，由 Zariski 拓扑中的相对开子集组成，且 \(\Omega \neq \emptyset\) 。若 \(\omega \in \Omega\) 且 \(x\) 是 \(\mathfrak{g}\) 的幂零元素，则映射 \(\tau : t \mapsto \omega \exp(t \operatorname{ad} x)\) （从 \(k\) 到 \(\mathrm{Aut}_e(\mathfrak{g})\) 的）是多项式的，故在 Zariski 拓扑中连续；于是 \(\tau(k)\) 是连通的；由于 \(\omega \in \tau(k)\) ，我们有 \(\tau(k) \subset \Omega\) 。这样，\(\Omega. (\exp \operatorname{ad} kx) \subset \Omega\) ，故 \(\Omega. \mathrm{Aut}_e(\mathfrak{g}) \subset \Omega\) 且 \(\Omega = \mathrm{Aut}_e(\mathfrak{g})\) 。这就证明 \(\mathrm{Aut}_e(\mathfrak{g})\) 是连通的。于是由 (i)，\(\mathrm{Aut}_0(\mathfrak{g})\) 是连通的。证毕。

我们将看到（§8 第 4 节 命题 6 的推论），\(\mathrm{Aut}_0(\mathfrak{g})\) 在 V 中对 Zariski 拓扑是闭的，且它是 \(\mathrm{Aut}(\mathfrak{g})\) 的单位元的连通分支。另一方面，\(\mathrm{Aut}_e(\mathfrak{g})\) 一般说来在 Zariski 拓扑中不是闭的。

* 假设 \((\mathfrak{g},\mathfrak{h})\) 是分裂的。群 \(\mathrm{Aut}_0(\mathfrak{g})\) 是群 \(G(k)\) ，即 \(k\) 上点的群，其中 \(G\) 是一个具有平凡中心的连通半单代数群（伴随群）。群 \(f(T_{Q})\) 等于 \(H(k)\) ，其中 \(H\) 是 \(G\) 的以 \(\mathfrak{h}\) 为李代数的嘉当子群。逆像 \(\tilde{H}\) （\(H\) 在泛覆盖 \(\tilde{G}\) （\(G\) 的，代数意义上的）中的逆像）以 \(T_{P}\) 为其 \(k\) -点群。\(\tilde{G}(k)\) 在 \(G(k) = \mathrm{Aut}_0(\mathfrak{g})\) 中的像是群 \(\mathrm{Aut}_e(\mathfrak{g})\) 。*

